\chapter{Esercizi: Strutture dati e tipi di dato astratti}
\label{chap:esercizi-strutture-dati-adt}

\begin{exercise}
\label{ex:fi-esercizi-strutture-dati-struttura-dati}
Che cosa si intende per \emph{data structure}?
\end{exercise}

Una \emph{struttura dati} specifica come gli elementi e le informazioni
necessarie a organizzarli sono rappresentati e memorizzati, così da consentire
l'esecuzione delle operazioni richieste (\cref{def:fi-struttura-dati}).

\begin{exercise}
\label{ex:fi-esercizi-strutture-dati-definizione-adt}
Che cosa si intende per \emph{Abstract Data Type} (ADT)?
\end{exercise}

Un \emph{tipo di dato astratto} descrive i valori e le operazioni di un
contenitore, specificandone il comportamento, le precondizioni e i casi di
errore senza imporre una rappresentazione interna (\cref{def:fi-adt}).

\begin{exercise}
\label{ex:fi-esercizi-strutture-dati-adt-struttura-differenza}
Qual è la differenza fondamentale tra ADT e \emph{data structure}?
\end{exercise}

L'ADT stabilisce \emph{quali} operazioni sono disponibili e \emph{quale}
risultato devono produrre; la struttura dati stabilisce \emph{come}
rappresentare gli elementi e realizzare tali operazioni
(\cref{sec:fi-adt-strutture}).

\begin{exercise}
\label{ex:fi-esercizi-strutture-dati-cosa-come}
Spiega la seguente frase: un ADT definisce \emph{cosa} un contenitore di dati
fa, mentre una struttura dati definisce \emph{come} viene implementato.
\end{exercise}

Per una coda, la specifica astratta impone che il primo elemento inserito sia il
primo rimosso. Questa proprietà non prescrive se usare un array circolare o una
lista concatenata: entrambe le strutture possono realizzare lo stesso
comportamento FIFO (\cref{sec:fi-adt-strutture,sec:fi-adt-coda}).

\begin{exercise}
\label{ex:fi-esercizi-strutture-dati-scelta-adt-struttura}
Perché la scelta di un ADT riguarda soprattutto l'\emph{efficacia} di una
soluzione? E perché la scelta di una \emph{data structure}, con cui implementare
un ADT, riguarda soprattutto l'\emph{efficienza} di una soluzione?
\end{exercise}

L'ADT viene scelto in base alle operazioni e alle proprietà richieste dal
problema: una scelta inadeguata non esprime correttamente il comportamento
desiderato. Una volta individuato l'ADT, la struttura dati determina soprattutto
i costi temporali e spaziali delle operazioni previste
(\cref{sec:fi-adt-strutture}).

\begin{exercise}
\label{ex:fi-esercizi-strutture-dati-stesso-adt-prestazioni}
Perché due implementazioni dello stesso ADT possono avere lo stesso
comportamento esterno ma \emph{performance} diverse?
\end{exercise}

La specifica astratta vincola i risultati osservabili, non il procedimento usato
per ottenerli. Una lista realizzata con array permette, per esempio, l'accesso
per indice in \(\Theta(1)\), mentre una lista concatenata richiede
l'attraversamento dei collegamenti (\cref{tab:fi-array-concatenata-costi}).

\begin{exercise}
\label{ex:fi-esercizi-strutture-dati-definizione-array}
Che cos'è un \emph{array}?
\end{exercise}

È una struttura dati composta da elementi di dimensione fissa memorizzati in
posizioni contigue, accessibili direttamente tramite indice
(\cref{def:fi-array}).

\begin{exercise}
\label{ex:fi-esercizi-strutture-dati-array-accesso-diretto}
Perché un array consente accesso diretto a un elemento tramite indice? Perché
tale accesso ha complessità \(O(1)\)?
\end{exercise}

L'indirizzo dell'elemento \(i\) si calcola come \(b+is\), dove \(b\) è
l'indirizzo iniziale e \(s\) la dimensione di ciascun elemento. Nel modello a
costo unitario il calcolo e l'accesso richiedono un numero costante di
operazioni, quindi \(\Theta(1)\) (\cref{sec:fi-strutture-array}).

\begin{exercise}
\label{ex:fi-esercizi-strutture-dati-statico-dinamico}
Qual è la differenza tra \emph{static array} e \emph{dynamic array}?
\end{exercise}

Un array statico ha una dimensione fissata; un array dinamico distingue il
numero di elementi presenti dalla capacità allocata e può ridimensionare
quest'ultima quando necessario (\cref{sec:fi-array-dinamici}).

\begin{exercise}
\label{ex:fi-esercizi-strutture-dati-size-capacity}
Spiega la differenza tra \emph{size} e \emph{capacity} in un \emph{dynamic
array}.
\end{exercise}

La \emph{size} \(n\) è il numero di elementi presenti; la \emph{capacity} \(C\)
è il numero di celle disponibili, con \(0\le n\le C\). Le \(C-n\) celle
inutilizzate sono già allocate ma non appartengono alla sequenza logica
(\cref{sec:fi-array-dinamici}).

\begin{exercise}
\label{ex:fi-esercizi-strutture-dati-append-ammortizzato}
Perché l'operazione di inserimento alla fine di un \emph{dynamic array} ha costo
ammortizzato \(O(1)\)?
\end{exercise}

Finché esiste una cella libera, l'inserimento costa \(\Theta(1)\). Quando
l'array è pieno, un ridimensionamento può copiare \(\Theta(n)\) elementi;
aumentando geometricamente la capacità, tuttavia, il costo complessivo di \(m\)
inserimenti è \(\Theta(m)\), dunque il costo ammortizzato è \(\Theta(1)\)
(\cref{sec:fi-array-dinamici}).

\begin{exercise}
\label{ex:fi-esercizi-strutture-dati-array-inserimento-intermedio}
Perché l'inserimento in posizione intermedia di un array dinamico può richiedere \(O(n)\)
tempo?
\end{exercise}

Per inserire all'indice \(i\) occorre spostare a destra i \(n-i\) elementi
successivi, dopo aver eventualmente aumentato la capacità. Il costo è \(O(n)\) e
può raggiungere \(\Theta(n)\) (\cref{sec:fi-array-dinamici}).

\begin{exercise}
\label{ex:fi-esercizi-strutture-dati-array-cancellazione-shift}
Perché la cancellazione da una posizione intermedia di un array può richiedere
uno \emph{shift} degli elementi?
\end{exercise}

Dopo la rimozione dell'elemento di indice \(i\), i \(n-i-1\) elementi seguenti
devono essere spostati a sinistra per mantenere contigua la sequenza logica. Il
costo degli spostamenti può essere \(\Theta(n)\) (\cref{sec:fi-array-dinamici}).

\begin{exercise}
\label{ex:fi-esercizi-strutture-dati-array-vantaggi-lista}
Quali sono i principali vantaggi di un array rispetto a una \emph{linked list}?
E i principali svantaggi?
\end{exercise}

L'array offre accesso diretto per indice in \(\Theta(1)\), buona località
spaziale e non richiede un riferimento al nodo successivo per ogni elemento. In
compenso, gli inserimenti e le rimozioni in testa o nel mezzo possono comportare
spostamenti \(\Theta(n)\) e, se dinamico, ridimensionamenti
(\cref{sec:fi-strutture-array,tab:fi-array-concatenata-costi}).

\begin{exercise}
\label{ex:fi-esercizi-strutture-dati-quando-usare-array}
In quali situazioni un array è una buona scelta implementativa per un ADT? E in
quali situazioni è una scelta poco efficiente?
\end{exercise}

È indicato quando prevalgono accessi e aggiornamenti per indice, oppure
inserimenti in fondo a un array dinamico. È meno conveniente quando sono
frequenti inserimenti o rimozioni in testa o in posizioni intermedie, che
richiedono spostamenti
(\cref{sec:fi-array-dinamici,tab:fi-array-concatenata-costi}).

\begin{exercise}
\label{ex:fi-esercizi-strutture-dati-linked-list-definizione}
Che cos'è una \emph{linked list}?
\end{exercise}

È una lista concatenata di nodi, ciascuno contenente un elemento e uno o più
riferimenti ad altri nodi. I nodi non devono essere contigui in memoria
(\cref{def:fi-lista-concatenata}).

\begin{exercise}
\label{ex:fi-esercizi-strutture-dati-nodo-singly}
Quali informazioni sono contenute in un \emph{node} di una \emph{singly linked
list}?
\end{exercise}

Ogni nodo contiene l'elemento memorizzato e il riferimento \texttt{next} al nodo
successivo. Nell'ultimo nodo, \texttt{next} è nullo
(\cref{sec:fi-liste-concatenate}).

\begin{exercise}
\label{ex:fi-esercizi-strutture-dati-head-tail}
Che cosa rappresentano i riferimenti \emph{head} e \emph{tail}?
\end{exercise}

\texttt{head} individua il primo nodo, mentre \texttt{tail}, se mantenuto,
individua l'ultimo. In una lista vuota entrambi sono nulli; aggiornare
\texttt{tail} permette l'inserimento in fondo in tempo costante
(\cref{sec:fi-liste-concatenate-operazioni}).

\begin{exercise}
\label{ex:fi-esercizi-strutture-dati-head-attraversamento}
Perché in una \emph{singly linked list} è sufficiente conoscere il riferimento
alla \emph{head} per attraversare tutta la struttura?
\end{exercise}

Partendo da \texttt{head} si segue ripetutamente il riferimento \texttt{next}
fino a raggiungere il nodo con riferimento successivo nullo. Ogni nodo permette
così di individuare il seguente (\cref{sec:fi-liste-concatenate}).

\begin{exercise}
\label{ex:fi-esercizi-strutture-dati-link-hopping}
Che cosa si intende per \emph{pointer hopping} o \emph{link hopping}?
\end{exercise}

È l'attraversamento di una lista concatenata seguendo, uno dopo l'altro, i
riferimenti che collegano i nodi. A differenza di un array, non consente di
calcolare direttamente l'indirizzo di un nodo dal suo indice
(\cref{sec:fi-liste-concatenate}).

\begin{exercise}
\label{ex:fi-esercizi-strutture-dati-linked-accesso-indice}
Perché l'accesso all'elemento in posizione \(i\) di una \emph{singly linked
list} ha complessità \(\Theta(i+1)\)?
\end{exercise}

Dalla testa occorre seguire \(i\) collegamenti per raggiungere il nodo
desiderato, contando anche l'accesso iniziale. Il costo è
\(\Theta(i+1)\): è costante per \(i=0\) e cresce linearmente con
\(i\) per \(i\ge1\)
(\cref{sec:fi-liste-concatenate}).

\begin{exercise}
\label{ex:fi-esercizi-strutture-dati-inserimento-head}
Perché l'inserimento in \emph{head} ha complessità \(O(1)\)?
\end{exercise}

È sufficiente collegare il nuovo nodo alla vecchia testa e aggiornare
\texttt{head}, senza attraversare o spostare altri elementi. Il costo è
\(\Theta(1)\) (\cref{sec:fi-liste-concatenate-operazioni}).

\begin{exercise}
\label{ex:fi-esercizi-strutture-dati-inserimento-tail-senza-riferimento}
Perché l'inserimento in \emph{tail} ha complessità \(O(n)\) se non si mantiene
un riferimento diretto alla \emph{tail}?
\end{exercise}

Per individuare l'ultimo nodo occorre partire da \texttt{head} e attraversare la
lista. La ricerca costa \(\Theta(n)\); con un riferimento \texttt{tail}
aggiornato, il collegamento del nuovo nodo costa invece \(\Theta(1)\)
(\cref{sec:fi-liste-concatenate-operazioni}).

\begin{exercise}
\label{ex:fi-esercizi-strutture-dati-varianti-linked-list}
Qual è la differenza tra \emph{singly linked list}, \emph{doubly linked list} e
\emph{circular linked list}?
\end{exercise}

Nella lista semplice ogni nodo rimanda al successivo; in quella doppia rimanda
anche al precedente, permettendo la navigazione nelle due direzioni. In una
lista circolare l'ultimo nodo rimanda al primo, anziché a un riferimento nullo.
Circolarità e doppio collegamento sono proprietà che possono anche coesistere
(\cref{sec:fi-liste-concatenate-operazioni}).

\begin{exercise}
\label{ex:fi-esercizi-strutture-dati-doubly-rimozione-nodo}
Perché una \emph{doubly linked list} facilita la rimozione di un nodo noto?
\end{exercise}

Il riferimento \texttt{prev} consente di raggiungere direttamente il
predecessore: conoscendo già il nodo, è possibile aggiornare i collegamenti dei
vicini in \(\Theta(1)\) tempo. Nella lista semplice può invece essere necessario
cercare il predecessore (\cref{sec:fi-liste-concatenate-operazioni}).

\begin{exercise}
\label{ex:fi-esercizi-strutture-dati-linked-vantaggi-svantaggi}
Quali sono i principali vantaggi e svantaggi delle \emph{linked lists} rispetto
agli array?
\end{exercise}

Permettono inserimenti e rimozioni in tempo costante quando sono già disponibili
i riferimenti necessari e non richiedono celle contigue. Richiedono però memoria
aggiuntiva per i collegamenti, hanno minore località spaziale e l'accesso per
indice comporta un attraversamento
(\cref{sec:fi-liste-concatenate,tab:fi-array-concatenata-costi}).

\begin{exercise}
\label{ex:fi-esercizi-strutture-dati-confronto-array-linked}
Confronta array e \emph{linked list} rispetto ad accesso per indice,
inserimento, cancellazione e uso della memoria.
\end{exercise}

L'array ha accesso per indice \(\Theta(1)\), ma inserimenti e cancellazioni nel
mezzo possono costare \(\Theta(n)\) per gli spostamenti. La lista concatenata ha
accesso per indice \(\Theta(i+1)\), ma inserisce o rimuove in \(\Theta(1)\)
quando sono noti i collegamenti da modificare. L'array conserva celle contigue,
eventualmente con capacità inutilizzata; la lista richiede riferimenti
aggiuntivi per nodo
(\cref{tab:fi-array-concatenata-costi,sec:fi-array-dinamici}).

\begin{exercise}
\label{ex:fi-esercizi-strutture-dati-stack-definizione}
Che cos'è uno \emph{stack ADT}?
\end{exercise}

È una pila, ossia un ADT che consente inserimento, rimozione e ispezione
soltanto alla cima (\cref{def:fi-adt-pila}).

\begin{exercise}
\label{ex:fi-esercizi-strutture-dati-stack-lifo}
Che cosa significa che uno \emph{stack} segue il principio LIFO?
\end{exercise}

\emph{Last In, First Out} significa che l'ultimo elemento inserito è il primo a
essere rimosso. Tutte le operazioni principali interessano la cima
(\cref{sec:fi-adt-pila,fig:fi-pila}).

\begin{exercise}
\label{ex:fi-esercizi-strutture-dati-stack-operazioni}
Descrivi il significato delle operazioni \texttt{push(e)}, \texttt{pop()},
\texttt{top()}, \texttt{is\_empty()} e \texttt{size()}.
\end{exercise}

\texttt{push(e)} inserisce \(e\) in cima; \texttt{pop()} rimuove e
restituisce l'elemento in cima; \texttt{top()} lo consulta senza rimuoverlo.
\texttt{is\_empty()}
verifica se la pila è vuota e \texttt{size()} restituisce il numero di elementi.
\texttt{pop()} e \texttt{top()} richiedono una pila non vuota
(\cref{sec:fi-adt-pila}).

\begin{exercise}
\label{ex:fi-esercizi-strutture-dati-stack-traccia}
Considera uno \emph{stack} inizialmente vuoto. Quali valori vengono restituiti
dalla seguente sequenza?
\[
\begin{gathered}
  \texttt{push(5)},\ \texttt{push(3)},\ \texttt{pop()},\ \texttt{push(8)},\\
  \texttt{top()},\ \texttt{pop()},\ \texttt{pop()}.
\end{gathered}
\]
\end{exercise}

Le operazioni che restituiscono un elemento producono, nell'ordine, \(3,8,8,5\):
dopo \texttt{push(5)} e \texttt{push(3)}, la prima \texttt{pop()} restituisce
\(3\); \texttt{top()} legge \(8\) senza rimuoverlo, e le ultime due
\texttt{pop()} restituiscono \(8\) e \(5\). La pila finale è vuota
(\cref{sec:fi-adt-pila}).

\begin{exercise}
\label{ex:fi-esercizi-strutture-dati-stack-array-topidx}
Come può essere implementato uno \emph{stack} usando un \emph{dynamic array} e
un indice \texttt{topIdx}?
\end{exercise}

La cima coincide con l'ultimo elemento presente. L'indice \texttt{topIdx} vale
inizialmente \(-1\); \texttt{push(e)} aggiunge l'elemento in fondo e lo
incrementa, mentre \texttt{pop()} restituisce l'ultima cella occupata e lo
decrementa. \texttt{top()} legge la cella senza modificarla
(\cref{sec:fi-adt-pila-implementazione-array}).

\clearpage

\begin{exercise}
\label{ex:fi-esercizi-strutture-dati-stack-push-resize}
Perché \texttt{push(e)} in uno \emph{stack array-based} può occasionalmente
richiedere \(O(n)\) tempo?
\end{exercise}

Quando l'array è pieno, occorre allocare una capacità maggiore e copiare gli
\(n\) elementi già presenti. Quella singola operazione può quindi costare
\(\Theta(n)\)
(\cref{sec:fi-array-dinamici,sec:fi-adt-pila-implementazione-array}).

\begin{exercise}
\label{ex:fi-esercizi-strutture-dati-stack-push-ammortizzato}
Perché, con analisi ammortizzata, \texttt{push(e)} in uno \emph{stack
array-based} può essere considerato \(O(1)\)?
\end{exercise}

Con capacità aumentata geometricamente, le copie complessive dovute ai
ridimensionamenti su \(m\) inserimenti costano \(O(m)\). Dividendo il costo
totale per \(m\), si ottiene un costo ammortizzato \(O(1)\) per
\texttt{push(e)}, non necessariamente per ogni singola chiamata
(\cref{sec:fi-array-dinamici}).

\begin{exercise}
\label{ex:fi-esercizi-strutture-dati-stack-linked-implementazione}
Come può essere implementato uno \emph{stack} usando una \emph{singly linked
list}?
\end{exercise}

Si fa coincidere la cima con \texttt{head}: \texttt{push(e)} inserisce un nodo
in testa, \texttt{pop()} rimuove la testa e \texttt{top()} ne legge l'elemento.
Tutte queste operazioni hanno costo \(\Theta(1)\)
(\cref{sec:fi-adt-pila,sec:fi-liste-concatenate-operazioni}).

\begin{exercise}
\label{ex:fi-esercizi-strutture-dati-stack-head-come-top}
Perché, in uno stack basato su lista concatenata, conviene usare
\emph{head} come \emph{top}?
\end{exercise}

Inserimento, rimozione e lettura in testa richiedono \(\Theta(1)\) e non
richiedono di conoscere il predecessore. Usando la coda di una lista semplice,
la rimozione dell'ultimo nodo richiederebbe invece una ricerca
(\cref{sec:fi-liste-concatenate-operazioni,sec:fi-adt-pila}).

\begin{exercise}
\label{ex:fi-esercizi-strutture-dati-stack-linked-costi}
Quali operazioni dello \emph{stack} sono \(O(1)\) in una \emph{linked-list-based
implementation}?
\end{exercise}

\texttt{push(e)}, \texttt{pop()}, \texttt{top()} e \texttt{is\_empty()} sono
\(\Theta(1)\) con la cima in \texttt{head}. Anche \texttt{size()} è
\(\Theta(1)\) se si mantiene un contatore aggiornato (\cref{sec:fi-adt-pila,sec:fi-liste-concatenate}).

\begin{exercise}
\label{ex:fi-esercizi-strutture-dati-stack-size-lineare}
Perché \texttt{size()} può essere \(O(n)\) in una \emph{linked list} se non si
mantiene una variabile aggiuntiva?
\end{exercise}

Senza un contatore, per conoscere la lunghezza è necessario attraversare tutti
gli \(n\) nodi. Il costo è \(\Theta(n)\); un contatore mantenuto dalle
operazioni di modifica lo riduce a \(\Theta(1)\) (\cref{sec:fi-liste-concatenate,sec:fi-adt-pila}).

\begin{exercise}
\label{ex:fi-esercizi-strutture-dati-stack-array-vs-linked}
Confronta l'implementazione \emph{array-based} e \emph{linked-list-based} di uno
stack.
\end{exercise}

Con array dinamico \texttt{push()} ha costo \(O(1)\) ammortizzato, ma può
richiedere \(\Theta(n)\) durante un ridimensionamento; con lista concatenata,
\texttt{push()} e \texttt{pop()} sono \(\Theta(1)\) senza copie. L'array offre
migliore località spaziale, mentre i nodi richiedono riferimenti aggiuntivi;
entrambe le soluzioni realizzano lo stesso ADT
(\cref{sec:fi-adt-pila,sec:fi-array-dinamici,sec:fi-liste-concatenate}).

\begin{exercise}
\label{ex:fi-esercizi-strutture-dati-queue-definizione}
Che cos'è una \emph{queue ADT}?
\end{exercise}

È una coda, ossia un ADT che inserisce elementi in fondo e li rimuove
dall'inizio, mantenendo l'ordine di arrivo (\cref{def:fi-adt-coda}).

\clearpage

\begin{exercise}
\label{ex:fi-esercizi-strutture-dati-queue-fifo}
Che cosa significa che una \emph{queue} segue il principio FIFO?
\end{exercise}

\emph{First In, First Out} significa che il primo elemento inserito è anche il
primo a essere rimosso, a differenza della politica LIFO della pila
(\cref{sec:fi-adt-coda,sec:fi-adt-pila}).

\begin{exercise}
\label{ex:fi-esercizi-strutture-dati-queue-operazioni}
Descrivi il significato delle operazioni \texttt{enqueue(e)},
\texttt{dequeue()}, \texttt{first()}, \texttt{is\_empty()} e \texttt{size()}.
\end{exercise}

\texttt{enqueue(e)} inserisce \(e\) in fondo; \texttt{dequeue()} rimuove e
restituisce il primo elemento; \texttt{first()} lo consulta senza rimuoverlo.
\texttt{is\_empty()} verifica se la coda è vuota e \texttt{size()} ne
restituisce la lunghezza. \texttt{dequeue()} e \texttt{first()} richiedono una
coda non vuota (\cref{sec:fi-adt-coda}).

\begin{exercise}
\label{ex:fi-esercizi-strutture-dati-queue-traccia}
Considera una \emph{queue} inizialmente vuota. Quali valori vengono restituiti
dalla seguente sequenza?
\[
\begin{gathered}
  \texttt{enqueue(A)},\ \texttt{enqueue(B)},\ \texttt{enqueue(C)},\\
  \texttt{dequeue()},\ \texttt{first()},\ \texttt{dequeue()}.
\end{gathered}
\]
\end{exercise}

I valori restituiti sono, nell'ordine, \(A,B,B\). La prima \texttt{dequeue()}
rimuove \(A\); \texttt{first()} legge \(B\) senza rimuoverlo e la seconda
\texttt{dequeue()} rimuove \(B\). Nella coda rimane \(C\)
(\cref{sec:fi-adt-coda}).

\begin{exercise}
\label{ex:fi-esercizi-strutture-dati-queue-circular-array}
Perché un \emph{circular array} è adatto a implementare una \emph{queue}?
\end{exercise}

Le posizioni liberate da \texttt{dequeue()} vengono riutilizzate mediante indici
calcolati modulo la capacità. In questo modo non occorre spostare tutti gli
elementi dopo una rimozione in testa
(\cref{sec:fi-coda-circolare,fig:fi-coda-circolare}).

\begin{exercise}
\label{ex:fi-esercizi-strutture-dati-queue-front-size-capacity}
In un \emph{circular array}, a che cosa servono \texttt{frontIdx}, \texttt{size}
e \texttt{capacity}?
\end{exercise}

\texttt{frontIdx} indica la posizione fisica del primo elemento; \texttt{size}
conta quelli presenti; \texttt{capacity} è il numero di celle allocate. Queste
informazioni permettono di distinguere l'ordine logico dagli indici dell'array
sottostante (\cref{sec:fi-coda-circolare}).

\begin{exercise}
\label{ex:fi-esercizi-strutture-dati-queue-back-formula}
Spiega la formula:
\[
  \texttt{backIdx}=(\texttt{frontIdx}+\texttt{size})\bmod\texttt{capacity}.
\]
\end{exercise}

Per una coda non piena e con capacità positiva, la somma \texttt{frontIdx} +
\texttt{size} individua la posizione immediatamente successiva all'ultimo
elemento nell'ordine logico. Il modulo riporta l'indice nell'intervallo delle
celle disponibili e permette il ritorno all'inizio dell'array
(\cref{sec:fi-coda-circolare}).

\begin{exercise}
\label{ex:fi-esercizi-strutture-dati-queue-dequeue-front}
Perché l'operazione \texttt{dequeue()} in un \emph{circular array} richiede
l'aggiornamento di \texttt{frontIdx}?
\end{exercise}

Rimosso il primo elemento, il successivo diventa la nuova testa della coda.
Occorre quindi porre
\(\texttt{frontIdx}\gets(\texttt{frontIdx}+1)\bmod\texttt{capacity}\) e
decrementare \texttt{size}, senza spostare gli elementi
(\cref{sec:fi-coda-circolare}).

\begin{exercise}
\label{ex:fi-esercizi-strutture-dati-queue-enqueue-resize}
Perché l'operazione \texttt{enqueue(e)} può richiedere \emph{resizing}?
\end{exercise}

Se \texttt{size} raggiunge \texttt{capacity}, non vi sono celle libere. In una
coda circolare ridimensionabile occorre allocare più spazio e ricopiare gli
elementi nell'ordine logico; il singolo ridimensionamento costa \(\Theta(n)\),
mentre l'inserimento resta \(O(1)\) ammortizzato (\cref{sec:fi-coda-circolare}).

\begin{exercise}
\label{ex:fi-esercizi-strutture-dati-queue-linked-implementazione}
Come può essere implementata una \emph{queue} usando una \emph{singly linked
list}?
\end{exercise}

Si mantiene \texttt{head} per il primo elemento e \texttt{tail} per l'ultimo:
\texttt{enqueue(e)} collega il nuovo nodo dopo \texttt{tail} e
\texttt{dequeue()} rimuove \texttt{head}. Occorre aggiornare entrambi i
riferimenti nei passaggi fra coda vuota e non vuota
(\cref{sec:fi-coda-concatenata}).

\begin{exercise}
\label{ex:fi-esercizi-strutture-dati-queue-head-tail-utilita}
Perché in una \emph{linked-list-based queue} è utile mantenere sia \emph{head}
sia \emph{tail}?
\end{exercise}

\texttt{head} consente di individuare il primo nodo per \texttt{dequeue()},
mentre \texttt{tail} permette \texttt{enqueue(e)} in \(\Theta(1)\). Senza
\texttt{tail}, l'inserimento in fondo richiede un attraversamento;
dal solo \texttt{tail} di una lista semplice non si può raggiungere
la testa (\cref{sec:fi-coda-concatenata}).

\begin{exercise}
\label{ex:fi-esercizi-strutture-dati-queue-senza-tail}
Che cosa accade alla complessità di \texttt{enqueue(e)} se non si mantiene un
riferimento alla \emph{tail}?
\end{exercise}

Per aggiungere in fondo bisogna percorrere la lista dalla testa fino all'ultimo
nodo. Il costo diventa \(\Theta(n)\), anziché \(\Theta(1)\) con \texttt{tail}
aggiornato (\cref{sec:fi-coda-concatenata}).

\begin{exercise}
\label{ex:fi-esercizi-strutture-dati-queue-array-vs-linked}
Confronta \emph{circular-array-based queue} e \emph{linked-list-based queue}.
\end{exercise}

Nell'array circolare \texttt{enqueue()} è \(O(1)\) ammortizzato e può richiedere
una copia durante il ridimensionamento, mentre \texttt{dequeue()} è
\(\Theta(1)\). Con lista concatenata e riferimenti \texttt{head}/\texttt{tail},
entrambe sono \(\Theta(1)\); la prima sfrutta celle contigue, la seconda
richiede un riferimento per nodo
(\cref{sec:fi-coda-circolare,sec:fi-coda-concatenata}).

\begin{exercise}
\label{ex:fi-esercizi-strutture-dati-deque-definizione}
Che cos'è una \emph{double-ended queue}, o \emph{deque}?
\end{exercise}

È una coda doppia che permette inserimento, rimozione e ispezione a entrambe le
estremità, senza prevedere accesso arbitrario agli elementi interni
(\cref{def:fi-adt-deque}).

\begin{exercise}
\label{ex:fi-esercizi-strutture-dati-deque-generalizza}
Perché una \emph{deque} è più generale sia di uno \emph{stack} sia di una
\emph{queue}?
\end{exercise}

Limitando tutte le operazioni a una sola estremità si ottiene una pila LIFO;
inserendo in fondo e rimuovendo all'inizio si ottiene una coda FIFO. La deque
permette entrambe le modalità e altre combinazioni (\cref{sec:fi-adt-deque}).

\begin{exercise}
\label{ex:fi-esercizi-strutture-dati-deque-operazioni}
Descrivi il significato delle operazioni \texttt{add\_first(e)},
\texttt{add\_last(e)}, \texttt{delete\_first()} e \texttt{delete\_last()}.
\end{exercise}

Le operazioni di inserimento aggiungono \(e\) rispettivamente all'inizio
(\texttt{add\_first(e)}) e alla fine (\texttt{add\_last(e)}).
Le operazioni di rimozione restituiscono l'elemento iniziale
(\texttt{delete\_first()}) o finale (\texttt{delete\_last()}), eliminandolo
dalla deque. Richiedono una deque non vuota (\cref{sec:fi-adt-deque}).

\clearpage

\begin{exercise}
\label{ex:fi-esercizi-strutture-dati-deque-traccia}
Considera una \emph{deque} inizialmente vuota. Mostra il contenuto dopo ogni
operazione:
\[
\begin{gathered}
  \texttt{add\_last(A)},\ \texttt{add\_last(B)},\\
  \texttt{add\_first(C)},\ \texttt{delete\_last()},\ \texttt{add\_first(D)}.
\end{gathered}
\]
\end{exercise}

Indicando gli elementi dall'inizio alla fine, gli stati successivi sono
\[
  [A],\quad [A,B],\quad [C,A,B],\quad [C,A],\quad [D,C,A].
\]
L'operazione \texttt{delete\_last()} rimuove e restituisce \(B\)
(\cref{sec:fi-adt-deque}).

\begin{exercise}
\label{ex:fi-esercizi-strutture-dati-deque-applicazioni}
In quali applicazioni può essere utile inserire e rimuovere elementi da entrambe
le estremità?
\end{exercise}

Una deque è utile quando occorre gestire attività da entrambe le estremità di
una sequenza oppure verificare un palindromo confrontando progressivamente il
primo e l'ultimo carattere (\cref{sec:fi-adt-deque}).

\begin{exercise}
\label{ex:fi-esercizi-strutture-dati-deque-doubly}
Perché una \emph{doubly linked list} è una possibile implementazione naturale
per una \emph{deque}?
\end{exercise}

Con riferimenti a testa e coda e collegamenti \texttt{prev}/\texttt{next}, si
possono individuare e aggiornare i nodi a entrambe le estremità in \(\Theta(1)\)
tempo, anche durante le rimozioni
(\cref{sec:fi-liste-concatenate-operazioni,sec:fi-adt-deque}).

\begin{exercise}
\label{ex:fi-esercizi-strutture-dati-deque-circular-array}
Perché anche un \emph{circular array} può essere usato per implementare una
\emph{deque}?
\end{exercise}

Gli indici circolari permettono di avanzare o arretrare la posizione iniziale e
di utilizzare celle libere su entrambi i lati logici. Non servono spostamenti
degli elementi; con crescita geometrica della capacità, inserimenti e rimozioni
hanno costo \(O(1)\) ammortizzato
(\cref{sec:fi-adt-deque,sec:fi-coda-circolare}).

\begin{exercise}
\label{ex:fi-esercizi-strutture-dati-list-adt-definizione}
Che cos'è una \emph{List ADT}?
\end{exercise}

È un ADT che rappresenta una sequenza di elementi, anche ripetuti, nella quale
la posizione è significativa e gli elementi possono essere letti, modificati,
inseriti e rimossi mediante indici (\cref{def:fi-adt-lista}).

\begin{exercise}
\label{ex:fi-esercizi-strutture-dati-list-vs-stack-queue}
In che cosa una \emph{list} differisce da uno \emph{stack} e da una
\emph{queue}?
\end{exercise}

La lista consente operazioni su posizioni arbitrarie tramite indici. Pila e coda
limitano invece l'accesso alle estremità, imponendo rispettivamente le politiche
LIFO e FIFO (\cref{sec:fi-adt-lista,sec:fi-adt-pila,sec:fi-adt-coda}).

\begin{exercise}
\label{ex:fi-esercizi-strutture-dati-list-ordine}
Perché in una \emph{list} l'ordine degli elementi è parte della struttura?
\end{exercise}

La posizione identifica ciascuna occorrenza: scambiare due elementi può cambiare
il risultato di \texttt{get(i)}, e inserimenti o rimozioni modificano gli indici
successivi. La lista rappresenta quindi una sequenza, non un insieme
(\cref{sec:fi-adt-lista,sec:fi-adt-insieme}).

\begin{exercise}
\label{ex:fi-esercizi-strutture-dati-list-operazioni}
Descrivi il significato delle operazioni \texttt{get(i)}, \texttt{set(i,e)},
\texttt{add(i,e)}, \texttt{append(e)} e \texttt{remove(i)}.
\end{exercise}

\texttt{get(i)} legge l'elemento \(i\); \texttt{set(i,e)} lo sostituisce con
\(e\); \texttt{add(i,e)} inserisce \(e\) alla posizione \(i\), facendo avanzare gli
indici degli elementi successivi. \texttt{append(e)} aggiunge in fondo
e \texttt{remove(i)} rimuove e
restituisce l'elemento di indice \(i\) (\cref{sec:fi-adt-lista}).

\begin{exercise}
\label{ex:fi-esercizi-strutture-dati-list-append-add}
Qual è la differenza tra \texttt{append(e)} e \texttt{add(i,e)}?
\end{exercise}

\texttt{append(e)} inserisce sempre in fondo, ossia nella posizione \(n\).
\texttt{add(i,e)} permette di scegliere una posizione \(0\le i\le n\) e fa
avanzare gli indici degli elementi successivi: l'inserimento in
fondo è il caso \(i=n\)
(\cref{sec:fi-adt-lista}).

\begin{exercise}
\label{ex:fi-esercizi-strutture-dati-list-traccia}
Considera la seguente sequenza su una \emph{list} inizialmente vuota:
\[
\begin{gathered}
  \texttt{add(0,A)},\ \texttt{append(B)},\ \texttt{append(D)},\\
  \texttt{add(2,C)},\ \texttt{remove(1)},\ \texttt{set(1,X)}.
\end{gathered}
\]
Qual è il contenuto finale della lista?
\end{exercise}

Gli stati sono
\[
\begin{aligned}
  [A]&\to[A,B]\to[A,B,D]\to[A,B,C,D]\\
     &\to[A,C,D]\to[A,X,D].
\end{aligned}
\]
\texttt{remove(1)} restituisce \(B\); \texttt{set(1,X)} sostituisce \(C\) con
\(X\). Il risultato finale è \([A,X,D]\) (\cref{sec:fi-adt-lista}).

\begin{exercise}
\label{ex:fi-esercizi-strutture-dati-list-array-get-set}
Perché in una \emph{array-based list} \texttt{get(i)} e \texttt{set(i,e)} hanno
complessità \(O(1)\)?
\end{exercise}

La posizione fisica dell'elemento si calcola direttamente dall'indice, senza
esaminare quelli precedenti. Lettura e aggiornamento richiedono dunque
\(\Theta(1)\) nel modello di costo adottato
(\cref{sec:fi-strutture-array,sec:fi-adt-lista}).

\begin{exercise}
\label{ex:fi-esercizi-strutture-dati-list-array-add-remove}
Perché in una \emph{array-based list} \texttt{add(i,e)} e \texttt{remove(i)}
possono avere complessità \(O(n)\)?
\end{exercise}

Un inserimento intermedio sposta a destra gli elementi successivi, mentre una
rimozione li sposta a sinistra. Il numero di spostamenti può essere
proporzionale a \(n\); per \texttt{add(i,e)} può inoltre essere necessario un
ridimensionamento (\cref{sec:fi-array-dinamici,sec:fi-adt-lista}).

\begin{exercise}
\label{ex:fi-esercizi-strutture-dati-list-linked-indice}
Perché in una \emph{linked-list-based list} l'accesso tramite indice non è
diretto?
\end{exercise}

I nodi non occupano posizioni contigue calcolabili dall'indice. Per trovare
l'elemento \(i\) occorre seguire i riferimenti a partire dalla testa, con costo
\(\Theta(i+1)\) (\cref{sec:fi-liste-concatenate,sec:fi-adt-lista}).

\begin{exercise}
\label{ex:fi-esercizi-strutture-dati-list-linked-convenienza}
Quando una \emph{linked-list-based list} può essere più conveniente di una
\emph{array-based list}?
\end{exercise}

Quando sono frequenti inserimenti e rimozioni vicino alla testa o presso nodi
già individuati, mentre sono rari gli accessi per indice. In questi casi si
possono aggiornare pochi collegamenti anziché spostare numerosi elementi
(\cref{tab:fi-array-concatenata-costi,sec:fi-adt-lista}).

\begin{exercise}
\label{ex:fi-esercizi-strutture-dati-list-doubly-positional}
Perché una \emph{doubly linked list} può essere utile per una
\emph{position-based list}?
\end{exercise}

Se una posizione è rappresentata dal riferimento a un nodo già noto, i
collegamenti \texttt{prev} e \texttt{next} permettono inserimenti e rimozioni
locali in \(\Theta(1)\). Raggiungere una posizione a partire dal solo indice
richiede invece ancora una ricerca
(\cref{sec:fi-liste-concatenate-operazioni,sec:fi-adt-lista}).

\begin{exercise}
\label{ex:fi-esercizi-strutture-dati-map-definizione}
Che cos'è un \emph{Map ADT}?
\end{exercise}

È un ADT che conserva associazioni tra chiavi distinte e valori. Ogni chiave
individua un solo valore, ma valori uguali possono essere associati a chiavi
diverse (\cref{def:fi-adt-mappa}).

\begin{exercise}
\label{ex:fi-esercizi-strutture-dati-map-key-value}
Che cosa sono una \emph{key} e una \emph{value}?
\end{exercise}

La \emph{chiave} identifica un'associazione e viene utilizzata per cercarla o
aggiornarla; il \emph{valore} è il dato associato a quella chiave
(\cref{sec:fi-adt-mappa}).

\begin{exercise}
\label{ex:fi-esercizi-strutture-dati-map-chiavi-uniche}
Perché in una \emph{map} le \emph{key} devono essere uniche?
\end{exercise}

Una chiave deve identificare senza ambiguità una sola associazione. Ripetere la
stessa chiave con un valore diverso non introduce una seconda associazione: ne
aggiorna il valore (\cref{def:fi-adt-mappa,sec:fi-adt-mappa}).

\begin{exercise}
\label{ex:fi-esercizi-strutture-dati-map-valori-ripetuti}
È possibile che due \emph{key} diverse siano associate allo stesso \emph{value}?
Fai un esempio.
\end{exercise}

Sì: \(\{1001\mapsto\text{Alice},\ 1003\mapsto\text{Alice}\}\) contiene due
chiavi distinte associate allo stesso valore. L'unicità è richiesta alle chiavi,
non ai valori (\cref{def:fi-adt-mappa}).

\begin{exercise}
\label{ex:fi-esercizi-strutture-dati-map-operazioni}
Descrivi il significato delle operazioni \texttt{get(k)}, \texttt{put(k,v)},
\texttt{remove(k)} e \texttt{contains(k)}.
\end{exercise}

\texttt{get(k)} restituisce il valore associato a \(k\); \texttt{put(k,v)}
inserisce o aggiorna l'associazione; \texttt{remove(k)} la elimina e
\texttt{contains(k)} verifica se la chiave è presente. Il comportamento di
\texttt{get()} e \texttt{remove()} su una chiave assente dipende
dall'interfaccia (\cref{sec:fi-adt-mappa}).

\begin{exercise}
\label{ex:fi-esercizi-strutture-dati-map-put-nuova}
Che cosa accade quando \texttt{put(k,v)} viene eseguita con una \emph{key} non
ancora presente?
\end{exercise}

La mappa inserisce la nuova associazione \(k\mapsto v\), aumentando di uno il
numero delle chiavi presenti (\cref{sec:fi-adt-mappa}).

\begin{exercise}
\label{ex:fi-esercizi-strutture-dati-map-put-esistente}
Che cosa accade quando \texttt{put(k,v)} viene eseguita con una \emph{key} già
presente?
\end{exercise}

Il valore precedentemente associato alla chiave \(k\) viene sostituito con
\(v\). Il numero di associazioni non cambia, perché la chiave rimane unica
(\cref{sec:fi-adt-mappa}).

\begin{exercise}
\label{ex:fi-esercizi-strutture-dati-map-get-esempio}
Considera la seguente \emph{map}:
\[
  \{1001\mapsto\text{Alice},\ 1002\mapsto\text{Bruno},\
    1003\mapsto\text{Alice}\}.
\]
Qual è il risultato di \texttt{get(1002)}?
\end{exercise}

Il risultato è \texttt{Bruno}, perché questa è la stringa associata alla chiave
\(1002\) (\cref{sec:fi-adt-mappa}).

\begin{exercise}
\label{ex:fi-esercizi-strutture-dati-map-alice-due-volte}
Nella mappa precedente, perché la presenza del valore \texttt{Alice} due volte
non viola la definizione di \emph{Map ADT}?
\end{exercise}

Le associazioni hanno chiavi distinte, \(1001\) e \(1003\). L'ADT impone
l'unicità delle chiavi, ma consente valori ripetuti (\cref{def:fi-adt-mappa}).

\begin{exercise}
\label{ex:fi-esercizi-strutture-dati-set-definizione}
Che cos'è un \emph{Set ADT}?
\end{exercise}

È un ADT che rappresenta una collezione di elementi distinti, senza associarvi
valori e senza attribuire significato all'ordine (\cref{def:fi-adt-insieme}).

\begin{exercise}
\label{ex:fi-esercizi-strutture-dati-set-vs-map}
In che cosa un \emph{set} differisce da una \emph{map}?
\end{exercise}

Un set contiene soltanto elementi distinti; una map contiene associazioni
\emph{chiave-valore}, con chiavi distinte ma valori eventualmente ripetuti
(\cref{sec:fi-adt-insieme,sec:fi-adt-mappa}).

\begin{exercise}
\label{ex:fi-esercizi-strutture-dati-set-senza-duplicati}
Perché un \emph{set} non consente duplicati?
\end{exercise}

L'appartenenza è una proprietà booleana: un elemento o appartiene all'insieme
oppure non vi appartiene. Ripeterlo non crea una nuova occorrenza, a differenza
di quanto accade in una lista (\cref{def:fi-adt-insieme,sec:fi-adt-lista}).

\begin{exercise}
\label{ex:fi-esercizi-strutture-dati-set-add-esistente}
Che cosa accade se si esegue \texttt{add(e)} su un elemento già presente nel
\emph{set}?
\end{exercise}

L'insieme rimane invariato, poiché \texttt{add(e)} inserisce l'elemento soltanto
se non è già presente (\cref{sec:fi-adt-insieme}).

\begin{exercise}
\label{ex:fi-esercizi-strutture-dati-set-operazioni}
Descrivi il significato delle operazioni \texttt{add(e)}, \texttt{remove(e)},
\texttt{contains(e)} e \texttt{size()}.
\end{exercise}

\texttt{add(e)} inserisce \(e\) se assente; \texttt{remove(e)} lo elimina;
\texttt{contains(e)} verifica l'appartenenza e \texttt{size()} restituisce il
numero di elementi distinti. Il comportamento di \texttt{remove(e)} su un
elemento assente è definito dall'interfaccia concreta
(\cref{sec:fi-adt-insieme}).

\begin{exercise}
\label{ex:fi-esercizi-strutture-dati-set-operazioni-insiemistiche}
Siano
\[
  A=\{1,2,3,4\},\qquad B=\{3,4,5\}.
\]
Calcola \texttt{union(B)}, \texttt{intersection(B)} e \texttt{difference(B)}
rispetto ad \(A\).
\end{exercise}

Le tre operazioni producono, rispettivamente,
\[
  A\cup B=\{1,2,3,4,5\},\qquad
  A\cap B=\{3,4\},\qquad
  A\setminus B=\{1,2\}.
\]
La differenza non è simmetrica: qui si conservano gli elementi presenti in \(A\)
ma assenti da \(B\) (\cref{sec:fi-adt-insieme}).

\begin{exercise}
\label{ex:fi-esercizi-strutture-dati-set-sottoinsieme}
Che cosa verifica l'operazione \texttt{is\_subset(T)}?
\end{exercise}

Verifica se ogni elemento dell'insieme su cui viene chiamata è presente anche in
\texttt{T}, cioè se vale \(A\subseteq T\). L'uguaglianza fra gli insiemi è
ammessa (\cref{sec:fi-adt-insieme}).

\begin{exercise}
\label{ex:fi-esercizi-strutture-dati-set-elimina-duplicati}
Perché un \emph{set} è utile per eliminare duplicati da una collezione?
\end{exercise}

Inserendo gli elementi della collezione in un insieme, ogni valore distinto
viene conservato una sola volta. Per esempio, da \([4,2,2,8,6,2]\) si ottiene
\(\{2,4,6,8\}\); l'ordine originale non è garantito dalla specifica dell'ADT
(\cref{sec:fi-adt-insieme}).
